\chapter{File formats and data}
\label{ch:formats}

\section{Products}

\begin{center}
\begin{tabularx}{\textwidth}{@{}lX@{}}
\toprule
File & Content \\
\midrule
\file{name.fbk.md} & the report: summary, analysis sections, findings,
  provenance, references (complete citations + paste-ready BibTeX) \\
\file{name.fbk.json} & the same content as data: \code{title}, \code{sections},
  \code{findings} (severity, message, rule id, suggested fix, refusals),
  \code{provenance}, plus a few program-specific keys \\
\file{structure.fbk.inp} & a generated ORCA input (pure ASCII) \\
\file{input.fix\_\textit{variant}.inp} & a corrected input from \menu{9} \\
\bottomrule
\end{tabularx}
\end{center}

Nothing carries timestamps; products are deterministic.

\section{Input JSON: cross-level records (menu 7)}

\begin{codeblock}
[
  {"label": "No.1",
   "occupation": "(5f)^1 ...",        # optional, informational
   "energies": {"HF": -1930.87692, "CCSD(T)": -1932.97951}},
  {"label": "No.2", "energies": {"HF": -1930.87392, "CCSD(T)": -1932.97745}}
]
\end{codeblock}

The top level may also be an object with a \code{records} list. All records
must carry the same level names; labels must be unique.

\section{Input JSON: crystal-field fit (menu 10)}

\begin{codeblock}
{
  "point_group": "C3",                  # C3 | Oh | C1/none
  "J": 8.0,                             # the |J M> manifold
  "levels": [0.0, 12.3, ...],           # one energy per sampled state
  "coefficients": [[[re, im], ...], ...],  # amplitudes on |J M>, M = m - J
  "declaration": {                      # optional: the five A5 items
    "convention": "...", "projection": "...", "units": "cm^-1",
    "z_axis": "...", "origin": "..."},
  "compare": {                          # optional: a second parameter set
    "label": "...",
    "parameters": {"2,0": 1879.0, "2,2": -43.0},
    "declaration": {"projection": "..."}}
}
\end{codeblock}

Parameter keys are \code{"k,q"} strings or two-element lists. Pass the
amplitudes themselves, not their complex conjugates (the conjugation trap,
Chapter~\ref{ch:menu10}).

\section{Input JSON: magnetic-doublet table (menu 16)}

\begin{codeblock}
{
  "system": "...",                       # optional label
  "reference": 0,                        # optional: index or label
  "doublets": [
    {"label": "ground", "g": [0.41, 0.44, 9.04], "energy": 0.0,
     "axis3": [0.0, 0.0, 1.0]},          # unit axis of the largest g value
    {"label": "1st excited", "g": [0.41, 0.44, 9.04], "theta3": 9.52}
  ]
}
\end{codeblock}

Per doublet, \code{theta3} (degrees) or \code{axis3} says where the largest
$g$ axis points; the reference doublet may omit both (its \code{theta3} is 0
by definition when the other angles are computed from axes, and it must carry
\code{axis3} then). \code{energy} is optional and printed when present.

\section{Input JSON: cross-structure mapping manifest (menu 17)}

\begin{codeblock}
{
  "structures": [                         # one entry per structure, same basis
    {"name": "r=1.094", "export": "a.json"},
    {"name": "r=2.600", "export": "b.json"}
  ],
  "tau": 0.5,                             # optional: one threshold for both criteria
  "selections": {"r=1.094": [4, 5]},      # optional: orbitals to make consistent
  "active": {"r=1.094": [4, 5, 6],        # optional: active-space overlap check
             "r=2.600": [4, 5, 6]}        # (per adjacent pair, in list order)
}
\end{codeblock}

Export paths are relative to the manifest; \code{selections} and \code{active}
name 0-based orbital indices per structure name.  Without \code{tau} the
threshold is read from the non-matchable curve (the source's Eq.~(8)).

\section{Input JSON: guess-transfer manifest (menu 18)}

\begin{codeblock}
{
  "structures": [                         # the neighbours, same basis/atom order
    {"name": "r=1.094", "export": "a.json"},
    {"name": "r=2.600", "export": "b.json"}
  ],
  "template": {                           # the target geometry
    "export": "target.json",              # its metric (S-Matrix)
    "mkl": "target.mkl"},                 # orca_2mkl <target> -mkl
  "delta": 1.0                            # optional neighbourhood radius (Angstrom)
}
\end{codeblock}

The output \code{<template stem>.fbk.mkl} converts with
\code{orca\_2mkl <template stem>.fbk -gbw} and is read with \code{!moread} +
\code{\%moinp}.  The mkl is a plain-text Molekel file (geometry in \AA,
verbatim \code{\$BASIS}, grouped coefficients, occupations); the toolkit's
reader/writer translates its measured p-shell row order to and from the
export convention.

\section{Input JSON: AOP rotation manifest (menu 47)}

\begin{codeblock}
{
  "reference": {                          # the active space to reproduce
    "export": "cas.json",                 # its orca_2json export
    "active": [4, 5, 6, 7, 8, 9]},        # 0-based active-orbital indices
  "target": {                             # the orbital set to start from
    "export": "scf.json",                 # same basis and atom order
    "mkl": "scf.mkl"},                    # orca_2mkl <scf> -mkl
  "closed": 4                             # the target's closed-shell orbital count
}
\end{codeblock}

The output \code{<target stem>.aop.fbk.mkl} converts with
\code{orca\_2mkl <target stem>.aop.fbk -gbw}.  The reference's active list
must have the same length as the target's active block; the containment
gate (Chapter~\ref{ch:menu47}) requires the reference to be represented in
the target's lowest \code{closed + active} orbitals.

\section{Input JSON: DM-AS selection manifest (menu 20)}

\begin{codeblock}
{
  "candidates": [                         # one entry per scanned space
    {"nel": 6, "norb": 8, "output": "cand_e6o8.out"}
  ],
  "reference": {"output": "dft.out"},     # the SCF block of that output is used
  "protocol": "gdm"                       # optional: gdm (default) | vgdm
}
\end{codeblock}

The reference may instead be a supplied value,
\code{\{"magnitude\_debye": 1.85, "source": "NIST (gas phase)"\}}, optionally
with a \code{vector} for the directional protocol.  Candidate outputs are
read at their single-root dipole block (State: 0); a state-averaged output is
refused.  Menu~\menu{19} writes this manifest (with the candidate list
pre-filled) next to the batch.

\section{The orca\_2json export and its request file (menus 12--15, 17, 21)}

\code{orca\_2json <base>.gbw} rewrites an orbital file into
\code{<base>.json}; which blocks it carries depends on the request file
\code{<base>.json.conf} placed next to the \code{.gbw}.  The requests the
menus use:

\begin{center}
\begin{tabularx}{\textwidth}{@{}>{\raggedright\arraybackslash}p{4.6cm} X@{}}
\toprule
Request & Blocks and measured conventions \\
\midrule
\code{\{"MOCoefficients": true\}} & Orbitals, occupations, energies; the base of
  every export. \\
\code{"1elIntegrals": ["H", "S", "T", "V"]} & \code{S-}/\code{H-}/\code{T}-/\code{V-Matrix};
  the cross-structure mapping (menu~\menu{17}) reads the kinetic matrix. \\
\code{"FockMatrix": ["J", "K"]} & Density-derived Fock blocks, one square
  matrix per spin; menu~\menu{21}.  Measured: the exported \code{K} block is
  the exchange contribution as it enters the Fock, so the source's
  $\tfrac12 K_{aa}$ equals $-\operatorname{diag}(CKC^{\mathsf T})_a =
  \sum_i (ai|ai)$; the exported \code{F} block equals $J+K$ and \emph{excludes}
  the core Hamiltonian ($H$ must be requested too for the full Fock). \\
\code{"2elIntegrals": ["MO\_IAJB"]},\newline
\code{"OrbWin": [i0, i1, a0, a1, 0, 0, 0, 0]} & Windowed two-electron
  entries \code{[i, j, a, b, value]} with \code{value} the chemist-notation
  $(ia|jb)$ and the internal indices first; menu~\menu{21} (\code{apcx}).  The
  input window is \emph{eight} integers (the second window all zeros for one
  window; the four-integer form is rejected by \code{orca\_2json}), inclusive
  0-based; the record written back into the JSON is the four-integer form. \\
\bottomrule
\end{tabularx}
\end{center}

Measured on ORCA 6.1.1: a \code{FockMatrix} request is answered from the run's
\code{<base>.densities} and \code{<base>.densitiesinfo} files, so the export
must be produced where the run wrote (or with the whole file family moved
together, names kept) -- a bare \code{.gbw} copy fails with ``NO densities are
available''.

\section{The ASS1ST round files (menus 22 and 23)}

A round is one CASSCF + FIC-NEVPT2 calculation with the density kept.  The
generated input (menu~\menu{22}, and menu~\menu{23} for the suggested space)
has the shape:

\begin{codeblock}
! <keywords> FIC-NEVPT2 KeepDens
%maxcore <MB>
%casscf
  nel <n>          # the active space of this round
  norb <m>
  mult <M>
  nroots <R>       # >1 makes it a state-averaged round
  PTSettings
    Density Unrelaxed
    NatOrbs true
  end
end
* xyz <charge> <M>
...
\end{codeblock}

Every piece is load-bearing (measured): \code{FIC-NEVPT2} because ORCA answers
the unrelaxed density only for the FIC ansatz; \code{KeepDens} because the
density is read from the run's \code{<base>.densities} sidecar (under
\code{Tdens-CASNEV.mult.<M>.root.<R>.p}, one entry per state); the
\code{PTSettings} block because \code{Density Unrelaxed} makes the density the
perturbation-theory one and \code{NatOrbs true} keeps the natural-orbital
basis.  The companion export request, written as \code{<base>.json.conf}:

\begin{lstlisting}[style=fbkcode]
{ "MOCoefficients": true, "1elIntegrals": ["S"], "Densities": ["all"] }
\end{lstlisting}

Round names chain: menu~\menu{22} writes \file{<stem>.r1.*}, and
menu~\menu{23}, fed \file{<stem>.rN.json}, writes \file{<stem>.r\{N+1\}.inp}
and its \code{.json.conf} beside the report.

\section{Input YAML: Judd-Ofelt dataset (menu 34)}

\begin{codeblock}
# one file per fit; comments and provenance lines are welcome
reference: "Babu et al., Physica B 279, 262 (2000) SET B"   # optional
host:
  refractive_index: 1.57          # a constant index, OR the Sellmeier form:
  # sellmeier:
  #   n0: 1.0
  #   terms:
  #     - [1.242781, 0.023833]    # A_i, B_i; B in um^2
transitions:
  - label: "7F6 <- 7F0"           # free text (printed as given)
    wavelength_nm: 2087           # OR energy_cm1; the wavelength also drives
                                  # the Sellmeier lookup when given
    f: 1.232e-6                   # measured absorption oscillator strength
    f_md: 0.0                     # optional: magnetic-dipole part (subtracted)
    j_low: 0                      # J of the lower level (2J+1 enters the fit)
    u2: 0.0                       # [U^(2)]^2, host-independent table value
    u4: 0.0
    u6: 0.1449
emission:                         # optional: radiative rates from the fit
  - label: "5D0 -> 7F2"
    energy_cm1: 16200
    f: 0.0                        # not used on this side
    j_low: 2                      # final level
    j_up: 0                       # emitting level (2J_up+1 enters A_ED)
    u2: 0.0032
    u4: 0.0
    u6: 0.0
\end{codeblock}

The fit itself uses the ``transitions`` entries only; at least four are
required (the standard deviation divides by $N-3$), and a $U^{(\lambda)}$
column that is zero everywhere means that $\Omega_\lambda$ is reported as
not determined.  The report is written as \file{<dataset>.jo.fbk.md}.

\section{Input YAML: pNMR run file (menu 35)}

\begin{codeblock}
structure: molecule.xyz        # resolved relative to the run file
center_atom: 1                 # 1-based index of the paramagnetic centre
temperature: 300               # K (drives the susceptibility and the <SS>)
susceptibility:
  from: tensor | g | zfs
  route: nonsymmetric          # or symmetric (the two constructions)
  # from: tensor
  tensor: [[cx_xx, cx_xy, cx_xz], [..], [..]]   # 3x3
  tensor_unit: "m3"            # or "1e-32 m3" (the literature unit)
  # from: g  (Kramers doublet unless spin is given)
  g_matrix: [[..], [..], [..]]
  spin: 0.5
  # from: zfs  (g + zero-field splitting)
  D_cm1: -85.5
  E_over_D: 0.00281
  # any of g_matrix / D_cm1 / E_over_D may instead be read from an output:
  orca_output: qdpt.out        # the effective-Hamiltonian g, the preferred ZFS
nuclei: all                    # or a list of 1-based indices (the centre excluded)
\end{codeblock}

The report is written as \file{<runfile>.pnmr.fbk.md}; every atom except the
centre is reported, or the listed subset.

\section{Structure and charge input (menus 2 and 11)}

Standard XYZ (atom count, comment, \code{element x y z}). For \menu{11}, the
interactive prompts take per-element charges (\code{O=-2,H=0.4}) and optional
radial moments \code{r2,r4,r6} in \AA$^k$. An element present in the structure
but missing from the charge list is refused, with the missing symbols named.

\section{Knowledge data (inside the package)}

\begin{center}
\begin{tabularx}{\textwidth}{@{}lX@{}}
\toprule
File & Content \\
\midrule
\file{knowledge/sources.bib} & the reference database (BibTeX). Every
  literature-evidence item must carry a bibkey that resolves here; loading
  fails otherwise \\
\file{knowledge/rules/*.yaml} & the rule base: id, kind (recipe/diagnosis),
  condition tree, action, refusals, evidence, confidence, severity (diagnosis
  only) \\
\file{knowledge/basis\_ecp.yaml} & the basis/ECP recommendation table \\
\file{knowledge/templates/*.j2} & the input templates (ORCA, OpenMolcas) \\
\file{knowledge/toolindex/tools.yaml} & the external-tool index (relation and
  status per entry) \\
\file{ui/menu.yaml} & the menu definition: number, title, handler (append-only
  numbers) \\
\bottomrule
\end{tabularx}
\end{center}

\section{Plot-ready CSV companions (the report's curves)}
\label{sec:formats-plotcsv}

Menus whose report is a curve or a table of plottable points additionally
write a \code{<output>.<menu>.fbk.csv} companion beside the report.  The
file is deliberately plain: one header line (column names with their
units, underscore-separated, e.g.\
\code{temperature\_K,rate\_per\_s,tau\_s}), then one data row per point,
comma separated, at ten significant digits, with no comment lines and no
styling information; a text field is quoted only when it would otherwise
break the row.  One deliberate extension (the shared spectrum writer, W0):
its \code{meta} option may prepend \code{\#\ key:\ value} comment lines
naming the columns, and with no metadata given its output is exactly the
comment-free companion above.  It drops straight into spreadsheets, Origin, pandas or
gnuplot, and the report beside it carries the context (fit window,
coefficients, boundary notes); the values are the engine's printed
numbers, so a figure made from the companion is regenerated from the same
output.  Covered today: menu~34 (the fitted transition table, with the
emission block's rows when present), menu~36 (the Orca\_Magrelax
$\tau(T)$ table), menu~37 (the primary block's stick spectrum), and
menu~41 ($S_{\mathrm{mag}}(T)$ on the levels route; the $-\Delta S(T, H)$
table on the Maxwell route); further curve-reporting menus join under the
same convention.

\section{The report as an interface}

The Markdown report is for humans; the JSON companion is the scripting
interface. Its stability follows the release policy: within a minor version
the section titles and finding fields are stable; new fields may be added.
Finding \code{rule\_id}s are stable identifiers you may key on.
